% ১৪.৫ Constant
\section{Constant}

\begin{learningbox}
এই টপিক শেষে শিক্ষার্থী constant-এর ধরন ও ব্যবহার লিখতে পারবে।
\end{learningbox}

\begin{definitionbox}{Constant}
Program execution চলাকালে যার মান পরিবর্তন হয় না, তাকে constant বলা হয়।
\end{definitionbox}

\begin{center}
\begin{tabular}{|p{0.25\textwidth}|p{0.35\textwidth}|p{0.25\textwidth}|}
\hline
\textbf{ধরন} & \textbf{উদাহরণ} & \textbf{নোট} \\
\hline
Integer constant & \texttt{10}, \texttt{-5} & পূর্ণসংখ্যা \\
\hline
Floating constant & \texttt{3.14} & দশমিকসহ সংখ্যা \\
\hline
Character constant & \texttt{'A'} & single quote \\
\hline
String constant & \texttt{"ICT"} & double quote \\
\hline
Symbolic constant & \texttt{\#define PI 3.1416} & নাম দিয়ে মান \\
\hline
\end{tabular}
\end{center}
